\documentclass[11pt]{article}
\usepackage{preamble}
\renewcommand{\answer}[1]{}

\begin{document}

\ladderheading{AP Statistics \textperiodcentered\ Unit 3 \textperiodcentered\ Topic 3.11 \textperiodcentered\ B: 2026-12-14 \textperiodcentered\ E: 2027-01-25}{Reminders and Practice Completion}

\focusbanner{TODAY'S PROBLEM}

In a district experiment, 360 volunteers are randomly assigned: 180 get daily reminders and 180 use a standard platform. By Friday, 126 reminder and 105 standard students finish optional practice. For students like these volunteers, a 95\% interval for $p_R-p_S$ is $(0.018,0.215)$.\par\smallskip \textbf{Iris:} ``The positive interval and random assignment support a causal increase of 1.8 to 21.5 points for students like these volunteers.''\par \textbf{Cole:} ``There is a 95\% chance reminders raise completion for all district students by at least 10 points.''\par
\begin{center}
\textbf{Experiment counts}\par\smallskip
\begin{tabular}{|l|c|c|c|}
\hline
\textbf{Version} & \textbf{Complete} & \textbf{Other} & \textbf{Total} \\ \hline
\textbf{Reminder} & 126 & 54 & 180 \\ \hline
\textbf{Standard} & 105 & 75 & 180 \\ \hline
\end{tabular}
\end{center}
\begin{firsttakebox}{FIRST TAKE --- before the video (5 min)}
Which report would you share with classmates? Give one reason from the study.
\workspace{0.9in}
\end{firsttakebox}

\begin{callout}{\IconBook\ MATCH EVERY CLAIM TO THE INTERVAL AND DESIGN (keep on the page)}{calloutgreen}
A confidence interval for $p_1-p_2$ gives plausible values for that population difference in the stated order.
If 0 is inside, no difference is plausible; if the whole interval is positive, the evidence supports $p_1>p_2$.
A claimed minimum effect $d$ is supported only when every interval value is at least $d$.
The confidence level is the long-run percent of intervals from repeated use of the method that capture the parameter,
not the probability that a fixed parameter or one computed interval is correct. Random assignment can support causation
for the experimental units; random sampling is what supports generalization to the sampled population.
\end{callout}

\newpage
\textbf{Finish after the video:}\par\smallskip

\dokbadge{1}~\textbf{(a)} State what $p_R-p_S$ represents, and identify whether $0$ and $0.10$ are contained in the reported interval.\par
\workspace{0.7in}

\dokbadge{2}~\textbf{(b)} Interpret the 95\% confidence interval in context. Then interpret the 95\% confidence level as the long-run performance of the interval procedure.\par
\workspace{1.4in}

\dokbadge{3}~\textbf{(c)} Choose the warranted conclusion. Use the positions of 0 and 0.10, distinguish random assignment from volunteer recruitment, and diagnose Cole's probability, minimum-effect, and population claims.\par
\workspace{2.0in}

\begin{sentenceframebox}\raggedright\textbf{Frame:} The interval supports a positive effect because \blank[1.8in], but it does not establish a 10-point minimum because \blank[2.0in].\end{sentenceframebox}

\begin{sentenceframebox}\raggedright\textbf{Frame:} Random assignment permits \blank[1.5in], whereas volunteer recruitment limits \blank[2.3in].\end{sentenceframebox}

\begin{summaryexitbox}{EXIT --- turn this in}
One thing the video changed about my first take: \hrulefill\par
\workspace{0.4in}
\end{summaryexitbox}

\end{document}
